% Source: MQ16a.html
\documentclass[11pt]{article}
\usepackage[margin=1in]{geometry}
\usepackage{amsmath}
\usepackage{amssymb}
\usepackage{siunitx}
\usepackage{enumitem}

\title{MQ16a: Sound}
\author{}
\date{}

\begin{document}
\maketitle

\section*{MQ16a: Sound}

\begin{enumerate}[leftmargin=*,label=\textbf{Question \arabic*.}]
\item \hfill \textit{1 pts}

A charge explodes at a depth of 1,342 m in the ocean. How much time would it take for the sound wave reach a helicopter hovering at a height of 920 m in the 7°C air above the water's surface?

Bulk modulus and density of sea water: 2.34 x 10$^{9}$ Pa and 1025 kg/m$^{3}$.

Ratio of heat capacities and molar mass of air: 1.4 and 28.8 g/mol.





\emph{t} = \rule{2cm}{0.4pt} s

\item \hfill \textit{1 pts}

A sound wave in an unknown fluid medium takes 10.3 ms to travel 10.6 m. What is the bulk modulus of the fluid in GPa, if the density is 1,027 kg/m$^{3}$? (1 GPa = 10$^{9}$ Pa)



\emph{B} = \rule{2cm}{0.4pt} GPa

\item \hfill \textit{1 pts}

A siren is designed to create a decibel level of 114 dB at a distance of 9 m. Calculate the sound intensity in W/m$^{2}$at the same distance if there are 2 sirens all sounding at that same location.

\item \hfill \textit{1 pts}

A pneumatic drill creates a decibel level of 122 dB at a distance of 9 m. Calculate the \textbf{sound intensity level }in dBat a distance of 45 m.

\item \hfill \textit{1 pts}

A rocket engine creates a decibel level of 122 dB at a distance of 20 m. What decibel level is created by a cluster of 4 such engines at a distance of 42 m?

\item \hfill \textit{1 pts}

A tuning fork sounds at 493 Hz. When a second tuning fork is struck at the same time, a beat frequency of 0.9 Hz is heard. If the second sound is lower in frequency than the first, what is its frequency?

\item \hfill \textit{1 pts}

Speaker 1 is at \emph{x} = 0 and speaker 2 is to the right and facing it at \emph{x} = 7 m. The speakers are connected to the same amplifier, which plays a continuous tone at 423 Hz. Find the \emph{x}-coordinate in meters of the first point of constructive interference that is right of center, if the temperature in the room is 16°C.

Take the speed of sound to be $v=(331\;\mathrm{m/s})\sqrt{\frac{T}{273\;\mathrm{K}}}$ .



\emph{x} = \rule{2cm}{0.4pt} m

\item \hfill \textit{1 pts}

Speaker 1 is at \emph{x} = 0 and speaker 2 is to the right and facing it at \emph{x} = 3 m. The speakers are connected to the same amplifier, which plays a continuous tone at 450 Hz. Find the \emph{x}-coordinate in meters of the first point of destructive interference that is left of center, if the temperature in the room is 18°C.



Take the speed of sound to be $v=(331\;\mathrm{m/s})\sqrt{\frac{T}{273\;\mathrm{K}}}$ .



\emph{x} = \rule{2cm}{0.4pt} m

\end{enumerate}

\end{document}
